\section{Conclusions}
\label{p13:sec:conclusions}

\noindent
A paper that has argued, throughout, that its object is one no single framework possesses, and that the truth of such an object lives not in any one voice but in the dialogue among them, cannot end in a single concluding voice without betraying its own claim. The conclusions are therefore plural, and they are plural by necessity and not by modesty. Each of the frameworks the paper has drawn upon is given here its own conclusion: what it can say about the generation of trust, said as firmly as it can be said, and what it cannot say, marked as firmly. The point of the exercise is not a survey but an enactment: the paper's form is meant to be its last argument, the demonstration, in the shape of its own ending, that the generation of trust is known only polyphonically, in the dialogue of irreducible positions none of which governs the rest. The discipline the series imposes on such an ending is exact, and the paper holds to it: each framework must speak to the edge of what it can warrant, hard and fallible and willing to assert, for a polyphony of voices that each hedge into vagueness is no polyphony but an evasion. Each voice below says something, and says where it must stop.

\subsection*{The jurisprudential conclusion}
\label{p13:sub:concl-jur}

The philosophy of law can say this: that the anchor of legal validity has migrated, across the century of its own reflection, from force toward trust, and that the law has recorded the migration in its own most technical doctrine, conceding at its edge, in promissory estoppel, that reliance and not compulsion is the true ground of a promise's binding force. It can say that force was never the cause of obligation but a substitute, now failing, for the production of expectation. It cannot say what trust \emph{is}, nor how, mechanically, it is produced; jurisprudence can locate the anchor and trace its movement, but the nature of the thing the anchor is moving toward lies outside its competence, and it must hand that question to others. Its limit is that it can say where validity rests without being able to say what, at that resting place, validity consists in.

\subsection*{The predictive-coding conclusion}
\label{p13:sub:concl-pc}

The predictive account can say this: that trust is the precision a generative model assigns to the prediction that a promise will be kept; that this dissolves the paradox of a present expectation directed at a non-existent future, since prediction is the brain's standing activity and requires no causation running backward from the future; that the expectation is produced through three coupled channels whose multiplicative coupling makes genuine trust robust where any single channel is fragile; and that the crisis of the symbol is the rational collapse of a population's precision-weight on a channel that forgery has made unreliable. It cannot say that trust simply \emph{is} a precision-weighting and nothing more; it offers a description at the computational and implementational level, and it must not be mistaken for the truth beneath the others. Its limit is the limit of every mechanistic account: it says how the thing is computed, and is silent on what the computing is, lived from within, and on whether it is just.

\subsection*{The geometric conclusion}
\label{p13:sub:concl-geo}

The geometric reading can say this: that a promise kept across time traces a closed loop, and that the sign of the holonomy the loop accumulates distinguishes the bond that deepens from the one that merely repeats and the one that drains; that the genuine, non-coercive promise is the adiabatic loop whose phase is its own, closed by no external pump, and the coerced bond the driven path whose phase is borrowed; and that this furnished, at last, the missing structure\,---\,a spectral gap\,---\,for a problem the series had left open. It cannot, in this paper, furnish the phase with a fully rigorous carrier; it has worked as a structural analogy, and the formal closure lies in another paper. Its limit is that it shows the shape of the distinction with precision while deferring the full formalisation of what bears the phase.

\subsection*{The quantum-dynamical conclusion}
\label{p13:sub:concl-quantum}

The quantum-dynamical reading can say this: that relational trust shares the formal structure of an entangled state, a property of correlations and not of nodes; that its failure under the crisis of the symbol is a systemic decoherence, the correlational shell preserved while the coherence drains; and that its transmission along the chain is the structural counterpart of entanglement swapping through an intermediary. It cannot say\,---\,and must say, loudly, that it does not say\,---\,that trust is in any literal sense a quantum phenomenon. It is a structural isomorphism and not a physical reduction, one more formal language for a structure and not the structure's secret physical truth. Its limit is the limit of every analogy: it illuminates by a likeness, and the likeness is not an identity.

\subsection*{The psychoanalytic conclusion}
\label{p13:sub:concl-psych}

The psychoanalytic reading can say this: that the production of trust circles a place that cannot be filled\,---\,that genuine trust, like the genuine poem, encircles the empty place of the \emph{das Ding} rather than filling it with an idol, and so requires no victim, where the forged trust fills the place with a positive object and must consume someone to do it; that the wager of trust is a venture made across a lack no evidence closes; and that the great Other, the guarantor that would certify the crossing, does not exist. It cannot translate the lack it speaks of into the positive terms the other frameworks want; its truths are truths about an absence, and they resist the cashing-out that would betray them. Its limit is that it speaks most truly where it can least be made to settle into a positive doctrine.

\subsection*{The political-economic conclusion}
\label{p13:sub:concl-pol}

The political-economic reading can say this: that the forms of trust-production are thrown up by the material conditions of their time and change as those conditions change; that centralised, force-based justice was the historically adequate form for a society of market strangers, and is failing now not through a decay of values but through a change in the material base; that the decentralisation of trust-production is a recomposition already under way, forced by that change; and that its emancipatory potential is real but not guaranteed, for the new channels are, under the present relations of production, being captured by capital toward a domination subtler than the sword. It cannot, from its own resources alone, settle whether the recomposition will emancipate or enslave; that depends on whether the relations of production are transformed, which is not a question the analysis of trust can answer by itself. Its limit is that it can diagnose the contradiction without being able to guarantee its resolution.

\subsection*{The Daoist conclusion}
\label{p13:sub:concl-dao}

The Daoist reading can say this: that the genuine bond is the one sustained by no external pump, closing upon itself by the completeness of what each freely brings; that to grasp at the guarantee is to lose the thing\,---\,that the one who would secure the bond by force forecloses the very deepening that only the unforced loop accumulates; and that the eternal the betrothal vow speaks of is founded upon no external warrant but here, in this heart, which is itself the ground\,---\,\zh{本乎此心，自成永恒}. It cannot say this in the language of mechanism or of proof; its truth is shown in a way of holding the thing, not demonstrated in a way of establishing it, and to press it into demonstration is to lose it as the grasping loses the bond. Its limit is the limit of the unsayable it points to: it can indicate the ground, and cannot, without contradiction, prove it.

\subsection*{Why there is no conclusion above the conclusions}
\label{p13:sub:concl-meta}

It remains to say why the paper does not now gather these into one\,---\,why there is no synthetic conclusion above the plural conclusions, no master voice to adjudicate among the frameworks and pronounce the truth they severally approach. The reason is the paper's deepest commitment, and it is the same reason that runs through the whole: there is no metalanguage. There is no position above the frameworks from which their several truths could be translated, without loss, into a single truth that contained them; the great Other who could occupy that position does not exist, and to write a synthetic conclusion would be to claim his place. The frameworks are irreducible epistemic positions on one object, and the object is known\,---\,this is the paper's form enacting the paper's content\,---\,only in the dialogue among them, in the mutual illumination that is itself a good cycle, a polyphony that generates understanding by the movement among voices rather than by their resolution into one. To conclude above the conclusions would be to settle what must be walked, to cash the dialogue into a doctrine, to commit, in the paper's last act, the very failure of settlement the paper has diagnosed throughout. The absence of a master conclusion is therefore not an incompleteness; it is the paper's final and formal assertion of its thesis. The knower of trust, like the truster, must forgo the guarantee\,---\,must hold the plural truths in their dialogue without the false comfort of a synthesis\,---\,and the paper, declining to synthesise, does in its form what it has asked, throughout, that the truster do in his life: it bears the absence of a guarantor, and binds itself to the plural truth anyway.

\subsection*{The three concentric circles}
\label{p13:sub:concl-circles}

One thing may be said, not as a synthesis above the frameworks but as an observation that runs beneath them all, and it returns the paper to the shape it announced in the Prelude. The question of trust was posed across three concentric scales: the intimate vow between two; the promise carried outward to a family; and the public institutions of contract and the courts. And the paper's finding, traced at each scale by each framework, is that the ground of trust is, at all three, the same ground. At the most public scale, the validity of law was found to rest, beneath the sword, on a reciprocal keeping of faith; at the most intimate, the bond was found to rest on a self-legislating relational divinity that needs no sword; and these are not two grounds but one, seen at two distances. What holds a society together at its most public and what holds two people together at their most intimate is the same thing: a trust generated not by force but by the free, self-binding, mutually relied-upon keeping of faith, across a gulf no guarantor spans. The three circles are concentric because they share a centre, and the centre is the relational divinity that descends wherever a free being binds itself, in relation, without recourse to a sword. The most intimate and the most public meet at that centre; and the betrothal vow on its red silk, which no court will enforce, is found to hold within it the ground on which even the law, at the last, is found to stand.
